\uuid{Rrg2}
\titre{Changements de variables}
\auteur{hanani ; mzari}
\organisation{exo7}
\licence{CC BY-NC-SA 4.0}
\datecreate{2020-01-20}
\niveau{L1}
\module{Analyse}
\chapitre{Calcul d'intégrales}
\sousChapitre{Autre}
\difficulte{3}
\theme{Calculs d'intégrales, changements de variables}
\exo7qcmid{lille-2.integrales.024}

\contenu{
\begin{question}
\qclassification{127.99}
\qtags{motcle=changements de variables}
Parmi les affirmations suivantes, cocher celles qui sont vraies :
\begin{answers}
\good{Le changement de variable $t=\cos x$ donne
$$\displaystyle \int _{\pi/3}^{\pi/2}\frac{\mathrm{d}x}{\sin x(1+\cos x)}=\int _0^{1/2}\frac{\mathrm{d}t}{(1-t)(1+t)^2}.$$}
\good{$\forall t\in \Rr\setminus \{-1,1\}$, $\displaystyle \frac{4}{(1-t)(1+t)^2}=\frac{1}{1-t}+\frac{1}{1+t}+\frac{2}{(1+t)^2}$.}
\bad{Une primitive de $\displaystyle \frac{1}{(1-t)(1+t)^2}$ sur $]-1,1[$ est $\displaystyle \ln \frac{1+t}{1-t}-\frac{2}{1+t}$.}
\bad{$\displaystyle \int _{\pi/3}^{\pi/2}\frac{\mathrm{d}x}{\sin x(1+\cos x)}=\ln 3+\frac{2}{3}$.}
\end{answers}
\vskip3mm
\begin{explanations}
Avec $t=\cos x\Rightarrow \mathrm{d}t=-\sin x\, \mathrm{d}x$, on a : $t(\pi/3)=1/2$, $t(\pi/2)=0$ et
$$\displaystyle \int _{\pi/3}^{\pi/2}\frac{\mathrm{d}\, x}{\sin x(1+\cos x)}=\int _{\pi/6}^{\pi/3}\frac{\sin x\, \mathrm{d}x}{\sin ^2x(1+\cos x)}=\int _{1/2}^{0}\frac{-\mathrm{d}\, t}{(1-t^2)(1+t)}.$$
Or, $\displaystyle \frac{4}{(1-t)(1+t)^2}=\frac{1}{1-t}+\frac{1}{1+t}+\frac{2}{(1+t)^2}$. Donc, en tenant compte de la linéarité,
$$\int _{\pi/3}^{\pi/2}\frac{\mathrm{d}\, x}{\sin x(1+\cos x)}=\frac{1}{4}\left[\ln \frac{1+t}{1-t}-\frac{2}{1+t}\right]_0^{1/2}=\frac{1}{4}\left(\ln 3+\frac{2}{3}\right).$$
\end{explanations}
\end{question}
}
